\subsection{Overview of GenericAgent}
\label{sec:overview}

\subsubsection{Design Principles}

GenericAgent (GA) is designed to maximize context information density, a principle that guides all architectural decisions. The effectiveness of LLM reasoning depends not on context length, but on the amount of task-relevant information each token conveys. Redundant or irrelevant tokens dilute attention and increase computational cost without improving performance.

We operationalize this principle by constraining all context content to satisfy three properties: \emph{completeness}, \emph{conciseness}, and \emph{naturalness}. Completeness ensures that all necessary information is present, avoiding unsupported inference. Conciseness removes redundancy and minimizes expression length while preserving meaning. Naturalness maintains standard language and interaction formats, preventing semantic distortion caused by over-compression.

This principle is systematically realized through four core technical dimensions. First, inspired by RISC processor design, GA exposes a minimal and orthogonal set of eight atomic tools whose schemas are injected into every LLM call, reducing prompt overhead while preserving expressiveness. Second, instead of loading all historical knowledge into the context, GA maintains a lightweight L1 index composed of one-line metadata entries, performing relevance selection before retrieving full memory content to prevent unbounded context growth. Third, repeated task patterns are distilled into reusable standard operating procedures (SOPs) and further compiled into executable Skills, compressing multi-step reasoning into single \texttt{code\_run} calls. Fourth, web content is transformed into compact Markdown through a two-stage DOM simplification pipeline, with subsequent interactions updating the context incrementally to avoid redundant transmission of full page content. Section~\ref{sec:core_design} presents detailed designs of these components.

\subsubsection{System Positioning}
GA is a general-purpose LLM agent designed for long-term autonomous operation, supporting two complementary execution modes that share a unified architecture. Interactive Mode processes explicit user requests, including code execution, file operations, web browsing, and information retrieval. Reflect Mode operates in the background, triggered by schedules or external events, performing system-level tasks such as monitoring, memory consolidation, and experience distillation without user intervention. Both modes are implemented on top of a shared execution kernel and can seamlessly switch across different LLM backends, including Claude, OpenAI, and Gemini.

\subsubsection{Runtime Flow}

The runtime of GA is organized around an iterative Agent Loop, as illustrated in Figure~\ref{fig:agent_loop}. At task initiation, global memory is loaded from persistent storage and merged with the system prompt to construct the initial session context. The user input is then submitted to the LLM. If the model generates tool calls, each call is dispatched to its corresponding handler, and the returned result is appended to the context as a new message. This process continues iteratively, forming a closed-loop interaction between the LLM and external tools.

The loop terminates when no further tool calls are produced. To ensure bounded execution, GA enforces a maximum of 40 reasoning turns, issuing warning signals every 7 turns and triggering a forced user confirmation via \texttt{ask\_user} at turn 35. After task completion, GA optionally invokes a memory distillation pipeline in which a dedicated sub-agent extracts useful knowledge and experience from the execution trace and stores them in long-term memory. Both Interactive Mode and Reflect Mode follow the same execution flow, differing only in how tasks are triggered and how outputs are consumed.

% \subsubsection{Design Principles}

% Every design decision in GA is guided by a single core principle: \textbf{maximizing context information density}.
% The reasoning quality of an LLM depends not on the length of its context, but on the density of task-relevant information carried by each token, superfluous tokens dilute attention and increase cost without improving outcomes.
% To enforce this principle, GA ensures that all content admitted into the context window satisfies three properties:

% \emph{completeness} (all information necessary to accomplish the task must be included, preventing the model from making unsupported guesses),
% \emph{conciseness} (within the bounds of completeness, content is expressed in the shortest possible form; redundant or irrelevant information is eliminated),
% and \emph{naturalness} (content must follow standard language and interaction formats, as excessive compression or fragmented instructions can distort meaning and break semantic coherence).

% This principle is realized across four core technical dimensions.
% \textbf{Atomic tool set}: inspired by RISC processor design, GA exposes only eight minimal, orthogonal tools whose schemas are injected into every LLM call, reducing per-turn prompt overhead and improving efficiency.
% \textbf{Hierarchical memory management}: rather than loading all accumulated knowledge at startup---which would cause unbounded prompt growth---GA maintains an L1 index of one-line metadata entries so the agent first determines whether a memory is relevant before incurring the cost of loading its full content.
% \textbf{Self-evolution}: repeated task patterns are distilled into reusable standard operating procedures (SOPs) and eventually into directly callable Skills, collapsing what previously required multiple reasoning rounds into a single \texttt{code\_run} call and freeing the token budget for novel problems.
% \textbf{Context-compressed browser control}: raw web pages are converted to compact Markdown through a two-stage DOM simplification pipeline before entering the context, and subsequent interactions update the context incrementally rather than retransmitting the full page.
% Detailed design and implementation of each dimension are presented in Section~\ref{sec:core_design}.

% \subsubsection{System Positioning}

% GenericAgent (GA) is a general-purpose LLM agent designed for long-term autonomous operation.
% It operates across two complementary modes.
% \textbf{Interactive Mode} handles real-time user requests, code execution, file operations, web browsing, and information retrieval, responding directly to explicit commands.
% \textbf{Reflect Mode} is triggered automatically by schedules or external events to perform background tasks such as system monitoring, memory consolidation, and experience distillation, without any user-initiated request.
% Both modes share a unified execution kernel and support flexible switching among mainstream LLM backends including Claude, OpenAI, and Gemini.

% \subsubsection{Runtime Flow}

% The GA runtime is centered on the Agent Loop, illustrated in Figure~\ref{fig:agent_loop}.
% At task initiation, global memory is loaded from persistent storage, merged with the system prompt, and used to construct the session-level context.
% The user input is submitted to the LLM; if the response contains tool calls, each call is dispatched to its corresponding handler and the result is appended as a new user message, driving the next reasoning round.
% If the response contains no tool calls, the task is considered complete and the loop exits.
% To prevent runaway execution, a hard cap of 40 reasoning turns is enforced; warnings are issued every 7 turns and a forced user confirmation via \texttt{ask\_user} is triggered at turn 35.
% Upon task completion, GA optionally runs a memory distillation pipeline in which a dedicated sub-agent extracts valuable facts and experience from the current execution trace and persists them to long-term memory.
% Both operating modes follow this identical flow, differing only in task sourcing and output routing.


\subsection{Core Technical Design}
\label{sec:core_design}

Guided by the principle of context information density maximization, GA's technical design spans four tightly coupled dimensions: memory organization, experience evolution, tool design, and browser control. These four dimensions are mutually reinforcing: the memory system governs the quality of knowledge admitted into the context; the self-evolution mechanism continuously compresses the token cost of execution; the minimal tool set controls the prompt overhead of tool definitions themselves; and browser control targets the high-noise scenario of web interaction with specialized page-level compression.


\subsubsection{Hierarchical Memory System}
\label{sec:memory}

The memory system constitutes the first line of defense for context information density optimization. GA employs a four-layer memory architecture that, through hierarchical distillation and taxonomy-based management, ensures only the highest-density knowledge enters the context window. The governing principle is to retain only the critical information that the model cannot derive through reasoning alone, thereby altering model behavior at minimal token cost.

\paragraph{L0: Meta-Rule Layer.}
The L0 layer stores the system's core behavioral norms and adheres to a strict minimalism principle. Each meta-rule is highly condensed---for example, the action-verification principle is expressed in a single sentence: ``Only information that has been empirically verified through actual execution may be written to long-term memory,'' thereby preventing hallucination contamination at the root. These rules are loaded into the System Prompt at every agent startup, establishing foundational behavioral guidance at an extremely low token cost.

\paragraph{L1: Index Layer.}
The L1 layer provides rapid localization of relevant memories through a memory taxonomy. It stores only metadata, file paths, brief descriptions, and applicability conditions, rather than full content, avoiding unnecessary loading of complete memory entries. Before executing a task, the agent first consults the L1 index and selectively loads only the memories genuinely needed into the context. Even as the memory store grows large, the context contains only task-relevant information, achieving precise retrieval with minimal token expenditure. In practice, the L1 index is maintained under a hard limit of 30 lines (fewer than 1{,}000 tokens), ensuring that the retrieval overhead remains negligible regardless of the total memory volume.

\paragraph{L2: Fact Store.}
The L2 layer stores verified environmental information, organized thematically and condensed through progressive summarization. Only incremental information that the model cannot infer on its own is retained---project directory structures, API endpoints, configuration parameters, and the like. Every fact is validated through tool execution before being committed, ensuring reliability. Topic-based organization prevents information fragmentation and improves retrieval efficiency.

\paragraph{L3: Task-Level SOP Layer.}
The L3 layer crystallizes long-term memory into reusable Skills: highly refined records of key steps and caveats distilled from multiple task executions. Only the portions where the model is prone to error or cannot easily reason---exception handling sequences, tool invocation orderings, and common failure avoidance strategies---are included. SOP refinement is carried out automatically through the Reflect mode, transforming raw experience into reusable knowledge that guides future tasks at minimal prompt cost.

\paragraph{Working Memory and Update Mechanisms.}
Working memory further optimizes context density by storing temporary information relevant to the current task. At each reasoning turn, an anchor prompt injects a working memory block containing the current turn number, stored key information from the session, and pointers to related SOPs. Upon task completion, a distillation step extracts valuable information and writes it to L2 or L3, after which working memory is cleared to prevent context pollution in subsequent sessions. Long-term memory updates employ two complementary strategies: precision editing, which modifies only the changed portions, and reflective distillation, which compresses accumulated experience. Together, these ensure the memory system is continuously compressed and optimized over time, steadily reducing prompt cost.


\subsubsection{Reflection-Driven Self-Evolution}
\label{sec:evolution}

The self-evolution mechanism is the core driver behind GA's continuous improvement in token efficiency. Its central insight is that experience should be progressively compressed from natural language into directly invocable code, enabling complex tasks to be completed with the fewest tokens possible. The system distills experience through the Reflect mode, consolidates high-density knowledge into SOPs and then into code, and finally registers the results as reusable Skills, realizing a self-driven evolutionary cycle.

\paragraph{Stage 1: Experience Distillation into SOPs.}
Through scheduled Reflect-mode sessions, the system autonomously analyzes task execution logs, focusing on the highest-density segments: effective strategies from successful tasks and recurring errors from failed ones. The agent extracts key steps and caveats from multi-turn tool calls and LLM reasoning traces, compressing lengthy execution trajectories into L3 SOPs. Only the portions where the model is error-prone or unable to reason independently are retained; routine steps and information that can be reconstructed are omitted, achieving information compression and knowledge distillation.

\paragraph{Stage 2: SOP Codification.}
For tasks that have been validated multiple times and exhibit a highly structured execution pattern, the system evaluates whether the corresponding SOP can be implemented as a Python script. For fixed-step, highly parameterized tasks---data transformation, batch file processing, periodic API calls---codification compresses what originally required multiple rounds of reasoning and tool invocation into a single \texttt{code\_run} execution, yielding an order-of-magnitude reduction in token consumption.

\paragraph{Stage 3: Code Registration as Skills.}
Codified scripts are registered as Skills in the L3 layer and the L1 index is updated accordingly. When the agent subsequently encounters a similar task, it can invoke the relevant Skill directly without re-deriving the solution. The system continuously monitors the effectiveness of SOPs and Skills, optimizing underperforming workflows---simplifying steps, reordering tool calls, and updating caveats---so that execution efficiency and information density improve over time.

These three stages form a closed loop: reflection distills experience, SOPs accumulate, codification lowers execution cost, and Skills support exploration of new tasks. Through this mechanism, GA achieves a steadily decreasing token cost for recurring task types as experience accumulates, while the freed token budget is redirected toward more complex tasks.


\subsubsection{Minimal Atomic Tool Set}
\label{sec:tools}

Tool definitions must be injected into the context on every LLM call, making tool design a direct determinant of context information density. To minimize the overhead introduced by tool schemas, GA adopts a minimal atomic tool set comprising only eight base tools, achieving complex task completion through principled composition. Table~\ref{tab:tools} summarizes the complete set.

\begin{table}[t]
\centering
\caption{The eight atomic tools in GenericAgent.}
\label{tab:tools}
\begin{tabular}{ll}
\toprule
\textbf{Tool} & \textbf{Description} \\
\midrule
\texttt{code\_run}            & Execute Python/Shell scripts with timeout control and abort signals. \\
\texttt{file\_read}           & Read files by line range with keyword-based location support. \\
\texttt{file\_write}          & Create or overwrite files with specified content. \\
\texttt{web\_scan}            & Extract the current browser page as simplified HTML, returning \\
                               & compressed Markdown and a tab list. \\
\texttt{web\_execute\_js}     & Inject and execute JavaScript in the user's browser, capturing \\
                               & return values and DOM changes automatically. \\
\texttt{ask\_user}            & Synchronous blocking request for user input; the Agent Loop \\
                               & immediately exits the current turn upon invocation. \\
\texttt{memory\_checkpoint}   & Persist key information from the current session to working memory \\
                               & for retrieval in subsequent reasoning turns. \\
\texttt{long\_term\_memory}   & Read from or write to cross-session long-term memory (L1--L3 layers), \\
                               & enabling knowledge accumulation and selective retrieval across tasks. \\
\bottomrule
\end{tabular}
\end{table}

GA's tool design adheres to an atomicity principle: each tool provides an irreducible primitive capability---code execution, file read/write, web content retrieval, browser-side execution, user interaction, and memory management. This minimal yet complete set covers every essential pathway for agent--environment interaction while avoiding the redundancy inherent in domain-specific tools.

The minimal tool set yields two direct benefits. First, it significantly reduces tool-description overhead in the context. In multi-turn interaction scenarios, tool definitions are repeatedly injected; fewer tools translate into a sustained reduction in cumulative token consumption across turns. Second, it lowers the model's decision complexity. A small, clearly bounded tool set reduces selection ambiguity, enabling the model to make more stable tool-call decisions and thereby reducing errors and retries.

In terms of expressiveness, GA does not rely on specialized tools to extend functionality but instead achieves complex operations through atomic-tool composition. File updates, for instance, follow a read-modify-write pattern; web data extraction is accomplished via content retrieval, structural analysis, and browser execution. This approach trades a small number of additional steps for a substantial gain in overall context efficiency. External capabilities such as MCP integrations are encapsulated as Skills that are invoked indirectly through the atomic tools and loaded into the context only on demand. Memory management is likewise implemented through sub-agents that operate the file-oriented tools, avoiding the introduction of additional tool types. Tool interfaces are kept minimal, retaining only the parameters strictly necessary, further reducing invocation complexity and token overhead.


\subsubsection{Context-Compressed Browser Control}
\label{sec:browser}

Browser interaction is one of the most severe sources of information redundancy in agent systems. Conventional approaches typically return the full web page content, including vast amounts of task-irrelevant structure, styling, and script information, resulting in extremely low context information density. GA redefines the core objective of browser control as minimum necessary information retrieval, systematically compressing page content, interaction traces, and execution paths to prevent low-value information from entering the context.

\paragraph{Content Layer: Simplified HTML.}
GA introduces a simplified HTML (simp HTML) mechanism that performs structural parsing and filtering of web pages, retaining only semantic information and key structural elements before converting them into a compact Markdown representation. The process operates in two stages: a first-stage JavaScript-side DOM analysis that classifies every node by relevance (e.g., primary content, secondary content, or covered/redundant), followed by a second-stage Python-side cleanup using BeautifulSoup that strips residual noise. Compared with raw HTML, this method eliminates styling, scripts, and redundant nesting, dramatically reducing token overhead while preserving informational completeness. Large lists are detected and automatically truncated to a representative sample (e.g., three items), and a smart budget-allocation mechanism ensures that token consumption stays within predefined limits. Hidden elements are expanded on demand, enabling the agent to access potentially critical information that would otherwise be lost due to page-visibility constraints.

\paragraph{Interaction Layer: Incremental DOM Updates.}
Rather than returning the full page after every operation, GA performs incremental updates based on DOM changes, extracting only the content that was added or modified by the most recent action. This strategy avoids the repeated transmission of large volumes of unchanged information during multi-step interactions, keeping the context size stable as the task progresses and significantly improving the efficiency of long-chain browser workflows.

\paragraph{Execution Layer: TMWebDriver.}
The TMWebDriver component provides stable browser control through a dual-channel communication architecture: a WebSocket channel (port 18765) for real-time bidirectional messaging and an HTTP long-polling channel (port 18766) for reliable delivery, both interfacing with a Chrome extension. Adaptive communication and multi-tab parallel execution reduce unnecessary interaction rounds. The system further leverages the Chrome DevTools Protocol (CDP) to access network requests directly, allowing certain tasks to bypass page parsing entirely and retrieve results at a lower information-transfer cost.

Taken together, GA transforms browser control from page-level information retrieval into task-relevant information extraction. Through the coordinated design of content compression, incremental updates, and execution-layer optimization, the system effectively maximizes context information density and interaction efficiency in web-based scenarios.
